% ============================================================================
%  SECCIÓN 6.1: ANÁLISIS DE LA INTEGRACIÓN CON SUPABASE Y POSTGRESQL
% ============================================================================

\section{An\'alisis de la integraci\'on con Supabase y PostgreSQL}

En los avances anteriores la aplicaci\'on \nombreApp{} conservaba la
informaci\'on en el propio dispositivo (persistencia local) o la manten\'ia
en memoria. Esa estrategia tiene una limitaci\'on fundamental: los datos
no est\'an disponibles en otros dispositivos y dependen del almacenamiento
de cada usuario. En este avance se incorpora un backend en la nube
mediante \textit{Supabase}, que provee una base de datos \textit{PostgreSQL}
centralizada, API, autenticaci\'on y almacenamiento de archivos.

\subsection{El problema de los datos aislados}

Para comprender la necesidad de una base de datos centralizada conviene
analizar una situaci\'on concreta. Suponga que un administrador de
\nombreApp{} registra un escenario deportivo en su dispositivo. Si la
informaci\'on se almacena \'unicamente en ese tel\'efono, otro usuario que
utilice un dispositivo distinto no podr\'a ver el nuevo escenario:

\begin{figure}[H]
\centering
  \figuraAPA{fig:problema-datos-aislados}{Problema de los datos aislados: cada dispositivo queda con su propia informaci\'on}{%
  \begin{tikzpicture}[
    node distance=0.5cm,
    caja/.style={rectangle, draw=black!60, rounded corners,
      minimum width=1.9cm, minimum height=0.6cm, align=center,
      font=\footnotesize},
    flecha/.style={-Stealth, thick}
  ]
    \node[caja] (a) {Dispositivo A};
    \node[caja, right=of a] (app1) {App};
    \node[caja, right=of app1] (srv) {Servidor\\ centralizado};
    \node[caja, right=of srv] (app2) {App};
    \node[caja, right=of app2] (b) {Dispositivo B};
    \draw[flecha] (a) -- (app1);
    \draw[flecha] (app1) -- (srv);
    \draw[flecha] (srv) -- (app2);
    \draw[flecha] (app2) -- (b);
  \end{tikzpicture}
  }{Elaboraci\'on propia en base a la gu\'ia del avance.}
\end{figure}

Sin un servidor, el usuario A crea un registro y el usuario B no tiene c\'omo
conocerlo. La soluci\'on es que toda la informaci\'on se consulte y se
escriba en un mismo lugar accesible por ambos dispositivos: el servidor
centralizado. All\'i aparece la necesidad de una base de datos compartida.

\subsection{Diferenciaci\'on entre frontend y backend}

La aplicaci\'on se divide en dos grandes componentes:

\begin{itemize}
  \item \textbf{Frontend:} la aplicaci\'on m\'ovil desarrollada con Expo
    (\textit{React Native}). Abarca las pantallas, los botones, los
    formularios, la navegaci\'on y la interacci\'on del usuario. Su
    responsabilidad es presentar la informaci\'on y capturar las acciones
    del usuario.

  \item \textbf{Backend:} la infraestructura que gestiona los datos y los
    servicios. En este proyecto corresponde a Supabase, que aporta la base
    de datos PostgreSQL, la API, la autenticaci\'on y el almacenamiento
    de im\'agenes. Su responsabilidad es administrar la informaci\'on de
    forma centralizada y segura.
\end{itemize}

La idea fundamental es: la aplicaci\'on presenta e interact\'ua con la
informaci\'on; Supabase permite gestionarla de forma centralizada y
compartida entre todos los usuarios.

\subsection{Arquitectura cliente--servidor con Supabase}

La arquitectura del avance se representa en la figura
\ref{fig:arquitectura-supabase}: la aplicaci\'on se comunica por Internet
con Supabase, y este gestiona la base de datos PostgreSQL (API
\textit{PostgREST}), la autenticaci\'on y el almacenamiento de archivos.

\begin{figure}[H]
\centering
  \figuraAPA{fig:arquitectura-supabase}{Arquitectura con Supabase como backend}{%
  \begin{tikzpicture}[
    node distance=0.7cm,
    caja/.style={rectangle, draw=black!60, rounded corners,
      minimum width=5.6cm, minimum height=0.8cm, align=center,
      font=\small},
    flecha/.style={-Stealth, thick}
  ]
    \node[caja] (app) {APLICACI\'ON M\'OVIL \\[2pt]
      \footnotesize Pantallas, formularios y l\'ogica};
    \node[caja, below=of app] (sup)
      {SUPABASE \\[2pt] \footnotesize Backend, API, autenticaci\'on y storage};
    \node[caja, below=of sup] (pg)
      {POSTGRESQL \\[2pt] \footnotesize Tablas, registros y relaciones};
    \draw[flecha] (app) -- node[right, font=\footnotesize]
      {Internet} (sup);
    \draw[flecha] (sup) -- (pg);
  \end{tikzpicture}
  }{Elaboraci\'on propia en base a la gu\'ia del avance.}
\end{figure}

De esta manera los datos dejan de depender del dispositivo y quedan
disponibles para todos los usuarios de la aplicaci\'on.

\subsection{Evoluci\'on del proyecto}

La integraci\'on de Supabase es la \'ultima etapa de la evoluci\'on de
los datos del proyecto:

\begin{enumerate}
  \item \textbf{PDF~4:} l\'ogica de negocio y estado en memoria.
  \item \textbf{PDF~5:} persistencia local (sesi\'on con AsyncStorage)
    y primeras consultas al servicio de Supabase.
  \item \textbf{PDF~6:} base de datos centralizada con Supabase y
    PostgreSQL, con operaciones CRUD conectadas a datos reales.
\end{enumerate}

El resultado es una aplicaci\'on conectada a una base de datos centralizada
en la nube, en la que las interfaces dejan de mostrar datos escritos
directamente en el c\'odigo y pasan a mostrar informaci\'on proveniente
del servidor.